% ========================================
% ANEXOS (no cuentan en la extensión del informe)
% ========================================

\anexo{REPRODUCIBILIDAD Y REPOSITORIO}
\label{an:repro}

\textbf{Repositorio:} \url{https://github.com/witman92/ml-leakage-auditor-capstone}.

Todas las tablas y figuras de este informe se generan en \arch{04_experiments/results/} y se reproducen con:

\begin{footnotesize}
\begin{verbatim}
pip install -r requirements.txt
python 04_experiments/leakage_demo.py   # ejemplo visual (Figura 1.2)
python run_all.py                       # fases 0-5, estadistica, figuras y pruebas
python run_all.py --rapido              # verificacion rapida (2 datasets x 2 familias)
streamlit run 06_web_app/app.py         # plataforma web
\end{verbatim}
\end{footnotesize}

\textbf{Estructura del repositorio.} Cada carpeta numerada corresponde a una fase del protocolo (Tabla~\ref{tab:etapas}):

\begin{footnotesize}
\begin{verbatim}
ml-leakage-auditor-capstone/
|-- run_all.py                  # reproduce todo con un comando
|-- requirements.txt            # versiones fijas (Python 3.13)
|-- 01_data/                    # Fase 0: carga de los 6 datasets de scikit-learn
|-- 02_pipelines/               # Fase 1: banco de 96 pares (controlled_pairs/),
|                               #   reparados (repaired_pairs/) y Conjunto A (llm_generated/)
|-- 03_auditor/                 # Fase 2: AST, verificador en ejecucion, refactorizador, CLI
|-- 04_experiments/             # Fases 3-4: benchmark, EDA, estadistica, demo, linea base
|   `-- results/                #   tablas (.csv), figuras (.png), manifiesto de ejecucion
|-- 05_skills/                  # Fase 5: dos Agent Skills autocontenidas
|-- 06_web_app/                 # Fase 5: plataforma web en Streamlit
|-- 07_report/                  # informe LaTeX, registro de decisiones, checklist
`-- 08_tests/                   # pruebas automaticas y 13 casos escritos a mano
\end{verbatim}
\end{footnotesize}

Cada script está organizado en bloques numerados (\texttt{\# \%\% PASO n}) que se pueden ejecutar uno a uno en VS Code o en Colab, y cada bloque explica de dónde vienen los datos, qué hace la función, qué métrica o gráfico produce y cómo interpretarlo. \arch{run_benchmark.py} escribe además \arch{run_manifest.json} con las versiones de las librerías y la semilla, y ordena todas las tablas de salida, de modo que dos ejecuciones producen archivos idénticos. La carpeta se abre en VS Code con \arch{capstone.code-workspace}, que trae una tarea por fase. Las decisiones y desviaciones están en \arch{07_report/decision_log.md}.
